% !TEX program = tectonic
% FICHE HOUSE — la donnée du dépôt au backtest, en un seul document. A4 portrait, 2 colonnes.
% Synthèse GRAPHIQUE des fiches DONNÉES + BACKTEST : 3 actes, ~11 boîtes, une figure par boîte,
% puces courtes, deux tailles de texte. Remplace FICHE_DONNEES_HOUSE et FICHE_BACKTEST_HOUSE.
% Tous les chiffres viennent du run gelé de Portefeuilles_House_Complet.ipynb (canal OCR activé).
% Figures : figures_house/ (extraites de ce même run, garanties cohérentes).
\documentclass[10pt,a4paper]{article}
\usepackage[margin=1.05cm]{geometry}
\usepackage[T1]{fontenc}
\usepackage[utf8]{inputenc}
\usepackage{lmodern}
\usepackage[french]{babel}
\usepackage{amsmath,amssymb}
\usepackage{booktabs}
\usepackage{multicol}
\usepackage{xcolor}
\usepackage{graphicx}
\usepackage{enumitem}
\usepackage{tikz}
\usepackage[most]{tcolorbox}
\usetikzlibrary{positioning}

\definecolor{bleu}{RGB}{20,77,144}
\definecolor{vert}{RGB}{20,110,60}
\definecolor{orange}{RGB}{196,110,0}
\definecolor{rouge}{RGB}{160,30,30}
\definecolor{grisf}{RGB}{110,110,110}
\definecolor{grisc}{RGB}{205,205,205}

\setlength{\columnsep}{5.5mm}
\setlength{\parindent}{0pt}
\setlength{\parskip}{0pt}
\pagestyle{empty}
\setlist[itemize]{leftmargin=3.4mm,topsep=0.8mm,itemsep=0.7mm,parsep=0pt,label=\textbullet}

% ── ruban : 11 pastilles, la courante pleine ────────────────────────────────
\newcommand{\ruban}[1]{%
  \begin{tikzpicture}[baseline=-0.6ex, x=1.9mm, y=1.9mm]
    \foreach \k in {1,...,11}{%
      \ifnum\k=#1\relax
        \fill[white] (\k,0) circle (0.82);
      \else
        \fill[white, opacity=0.42] (\k,0) circle (0.52);
      \fi}
  \end{tikzpicture}}

\newtcolorbox{etapebox}[3]{%
  enhanced, breakable=false, colback=#2!3, colframe=#2!55, boxrule=0.5pt,
  arc=2pt, left=2.2mm, right=2.2mm, top=1.6mm, bottom=1.6mm,
  toptitle=0.6mm, bottomtitle=0.6mm, boxsep=0pt,
  title={\begin{minipage}[c]{0.60\linewidth}\raggedright\footnotesize\bfseries
           \textcolor{white}{#3}\end{minipage}\hfill
         \begin{minipage}[c]{0.33\linewidth}\raggedleft\ruban{#1}\end{minipage}},
  colbacktitle=#2, coltitle=white, fonttitle=\footnotesize\bfseries,
  before skip=1.7mm, after skip=0mm}

\newcommand{\reponse}[1]{{\footnotesize #1\par}\vspace{1mm}}
\newcommand{\phare}[3]{%
  \begin{center}\vspace{-0.4mm}
  {\Large\bfseries\color{#3} #1}\\[0.3mm]{\footnotesize #2}
  \end{center}\vspace{0.5mm}}
\newcommand{\pts}[1]{{\footnotesize\begin{itemize}#1\end{itemize}}\vspace{0.6mm}}
\newcommand{\vis}[1]{\vspace{0.6mm}\begin{center}\includegraphics[width=0.95\linewidth]{figures_house/#1}\end{center}\vspace{0.3mm}}
\newcommand{\visb}[2]{\vspace{0.6mm}\begin{center}\includegraphics[width=#2\linewidth]{figures_house/#1}\end{center}\vspace{0.3mm}}
\newcommand{\src}[1]{{\scriptsize\color{grisf} #1\par}}
\newcommand{\acte}[2]{%
  \vspace{1mm}{\footnotesize\bfseries\color{#2} #1}\vspace{-1.3mm}
  \par{\color{#2}\rule{\linewidth}{1pt}}\vspace{0.5mm}}
\newlength{\colw}\setlength{\colw}{0.4865\textwidth}
\newcommand{\rangee}[2]{%
  \noindent\begin{minipage}[t]{\colw}\vspace{0pt}#1\end{minipage}%
  \hfill\begin{minipage}[t]{\colw}\vspace{0pt}#2\end{minipage}\par}

\begin{document}

% ═══════════════ CHAPEAU ═══════════════
\begin{center}
{\large\bfseries Du dépôt au backtest --- la donnée House de bout en bout}\\[0.2mm]
{\scriptsize ce que les élus déposent $\rightarrow$ ce qu'on en reconstruit $\rightarrow$ ce que donne une
stratégie qui les copie $\cdot$ House 2014-2026 $\cdot$ notebook \texttt{Portefeuilles\_House\_Complet}}
\end{center}
\vspace{0.9mm}

\begin{tcolorbox}[enhanced, colback=white, colframe=grisc, boxrule=0.5pt, arc=2pt,
                  left=3mm, right=3mm, top=1.1mm, bottom=1.1mm]
{\footnotesize\bfseries\color{bleu} ACTE I --- la donnée}\hfill
{\footnotesize\bfseries\color{vert} ACTE II --- la reconstruction}\hfill
{\footnotesize\bfseries\color{orange} ACTE III --- la stratégie}\\[0.8mm]
\begin{center}
\begin{tikzpicture}[x=1.02cm, y=0.78cm, font=\scriptsize]
\foreach \k/\lab/\c [count=\i from 0] in {%
  1/pop./bleu, 2/att.\,vs\,obs./bleu, 3/invent./bleu, 4/tri/bleu,
  5/valoriser/vert, 6/reconstr./vert, 7/vérifier/vert,
  8/règle/orange, 9/réglage/orange, 10/{combien K}/orange, 11/portée/orange}{%
  \ifnum\k<11 \draw[grisc, line width=0.9pt] ({\i*1.34+0.26}, 0) -- ({\i*1.34+1.08}, 0);\fi
  \fill[\c] ({\i*1.34}, 0) circle (0.19);
  \node[text=white, font=\tiny\bfseries] at ({\i*1.34}, 0) {\k};
  \node[text=\c, font=\tiny, anchor=north, align=center] at ({\i*1.34}, -0.22) {\lab};}
\end{tikzpicture}
\end{center}
\vspace{0.3mm}
{\footnotesize\textbf{Les mots, une fois pour toutes.} Un \textbf{membre} dépose des \textbf{documents},
faits de \textbf{lignes} : une \textbf{position} détenue (\emph{ce que je possède}) ou une
\textbf{transaction} datée (\emph{ce que j'ai échangé}). On \textbf{reconstruit} le portefeuille de chacun,
puis une stratégie le \textbf{copie}.}
\end{tcolorbox}

\vspace{1mm}
% ═══════════════ ACTE I — LA DONNÉE ═══════════════
\acte{ACTE I --- la donnée : qui dépose quoi, et ce qu'on en garde}{bleu}
\rangee{\begin{etapebox}{1}{bleu}{1 $\cdot$ LA POPULATION --- qui est là, qui trade}
\reponse{900 membres, 4 trajectoires de mandat (haut) --- 408 tradent des actions cotées, vus par des
canaux différents (bas).}
\phare{408 sur 900}{membres qui tradent des actions cotées}{bleu}
\vis{fig_10_quadrant.png}
\vis{fig_10_trois_zones.png}
\src{86 membres ne sont vus \emph{que} par leurs rapports annuels ; 17 de plus par des documents scannés
($391 + 17 = 408$). Source : cellules \og qui trade \fg{} et \og validation croisée \fg{}.}
\end{etapebox}}{\begin{etapebox}{2}{bleu}{2 $\cdot$ A-T-ON CE QUI EST DÛ ? --- attendu vs observé}
\reponse{Un \textbf{arrivé} doit une photo d'entrée, un \textbf{partant} une photo de sortie (barres pleines
$=$ observé, pointillés $=$ attendu). Et chacun, un rapport chaque année.}
\phare{85 \%}{des années de mandat couvertes par un rapport (médiane)}{bleu}
\vis{fig_10_etape_arrivee.png}
\vis{fig_10_etape_sortie.png}
\src{une partie des manques est réelle (rien déposé), une partie vient de documents pas encore exploités ---
séparé membre par membre dans le notebook. Source : cellule \og tableau central \fg{}.}
\end{etapebox}}
\vspace{2mm}
\rangee{\begin{etapebox}{3}{bleu}{3 $\cdot$ L'INVENTAIRE --- tout ce qui est déposé}
\reponse{Premier travail : réunir tous les documents et compter ce qu'ils contiennent, sans rien écarter.}
\phare{609\,886}{lignes de patrimoine, dans 14\,887 documents}{bleu}
\vis{fig_10_inventaire.png}
\pts{
\item 4\,447 rapports annuels, 934 photos de candidat, 378 de sortie, 376 d'entrée --- ramenées à
325 photos d'entrée retenues, une par membre.
\item \textbf{Un tiers vient de documents scannés}, relus automatiquement : 188\,970 lignes, dont 81\,535
avec un symbole boursier. Chaque ligne garde la trace de son origine.
}
\src{source : cellules \og inventaire documentaire \fg{} et \og bouclage \fg{} (18 membres sans aucun document).}
\end{etapebox}}{\begin{etapebox}{4}{bleu}{4 $\cdot$ LE TRI --- action ou pas action}
\reponse{Avant tout prix, chaque ligne reçoit une classe : action cotée, fonds, monétaire, ou code de
formulaire. Le moteur ne verra que les actions et ETF.}
\phare{231\,426 sur 609\,886}{lignes qui sont vraiment des actions ou des ETF}{bleu}
\vis{fig_10_paysage.png}
\pts{
\item La majorité de ce qui est déclaré n'a pas de ticker : immobilier, sociétés non cotées, comptes.
\item En valeur cumulée : 87,4 Md\$ sans ticker contre 12,8 Md\$ d'actions. Le patrimoine \emph{réel}, au
dernier document de chaque membre, vaut \textbf{8,9 Md\$}.
}
\src{source : cellules \og classification \fg{} et \og paysage déclaré \fg{}. Valeurs = milieu de fourchette ; le cumul recompte un actif à chaque rapport, ce n'est pas un patrimoine.}
\end{etapebox}}
\vspace{2.6mm}

% ═══════════════ ACTE II — LA RECONSTRUCTION ═══════════════
\acte{ACTE II --- la reconstruction : de la ligne déclarée au portefeuille vérifié}{vert}
\rangee{\begin{etapebox}{5}{vert}{5 $\cdot$ VALORISER --- les transactions retenues}
\reponse{Une transaction n'entre au moteur que si l'on connaît le prix de son titre le jour dit. Trois
sources s'additionnent, chacune datée.}
\phare{129\,538}{transactions datées et valorisées}{vert}
\vis{fig_10_delais.png}
\pts{
\item \textbf{107\,976} déclarations rapides (publiques en 27 jours) $+$ \textbf{14\,084} des rapports
annuels (publiées 373 jours après) $+$ \textbf{7\,478} des documents scannés.
\item Deux horloges pour une même transaction : la rapide est un \emph{signal}, l'annuelle un
\emph{enregistrement}. Chaque transaction garde sa provenance.
}
\src{source : cellules \og livre de comptes O/A/T \fg{} et \og délais \fg{}. Chaque ligne écartée a une raison affichée.}
\end{etapebox}}{\begin{etapebox}{6}{vert}{6 $\cdot$ RECONSTRUIRE --- la photo devient un portefeuille}
\reponse{Une photo donne des fourchettes de valeur, pas des quantités. On la convertit en parts détenues
au prix du jour de son dépôt, puis on applique les transactions.}
\phare{243}{portefeuilles de départ construits}{vert}
\vis{fig_10_entonnoir_h0.png}
\begin{center}\footnotesize
$h_0 = \dfrac{\tilde v}{P(\text{jour du dépôt})}$ \quad achat $h \mathrel{+}= \dfrac{\tilde v}{P}$ \quad
vente $h \mathrel{-}= \min(q,h)$
\end{center}
\vspace{0.4mm}
\pts{
\item Prix du jour du \textbf{dépôt}, jamais antérieur : rien ne regarde le futur. 84,5 \% de la valeur en
actions des photos devient des parts.
\item Sans photo, un membre part de zéro : son point de départ n'est pas nul, il est inconnu.
}
\src{source : cellules \og $h_0$ en parts \fg{} et \og livre de comptes photos \fg{}.}
\end{etapebox}}
\vspace{2mm}
\rangee{\begin{etapebox}{7}{vert}{7 $\cdot$ VÉRIFIER --- les rapports annuels jugent}
\reponse{Le rapport annuel dit ce que l'élu possédait vraiment. On confronte notre reconstruction à cette
déclaration, position par position, à la date du dépôt.}
\phare{9 \% $\rightarrow$ 57 \%}{des positions déclarées retrouvées, sur la photo médiane}{vert}
\vis{fig_10_recouvrement.png}
\pts{
\item Avec les seules déclarations rapides, la reconstruction ne retrouve que 9 \% des positions déclarées.
Avec le \textbf{dossier complet} (photo d'entrée $+$ rapports annuels), \textbf{57 \%} --- six fois plus.
\item Sans perdre en justesse : la précision passe de 55 à 56 \%, la valeur de 44 à 47 \%. Médianes sur
1\,557 et 1\,878 photos. Même exercice sur 136 photos de sortie : 60 \% retrouvé.
}
\src{source : cellules \og confrontation \fg{} et \og agrégats \fg{}. C'est le juge indépendant de tout ce qui suit.}
\end{etapebox}}{\begin{tcolorbox}[enhanced, breakable=false, colback=grisc!12, colframe=grisc, boxrule=0.5pt,
                  arc=2pt, left=2.2mm, right=2.2mm, top=1.6mm, bottom=1.6mm,
                  toptitle=0.6mm, bottomtitle=0.6mm, boxsep=0pt, before skip=1.7mm, after skip=0mm,
                  colbacktitle=grisf, coltitle=white, fonttitle=\footnotesize\bfseries,
                  title={CHARNIÈRE --- ce qu'on a en main pour la suite}]
\reponse{Fin de la reconstruction. On dispose maintenant, pour chaque membre qui trade, d'un
\textbf{portefeuille daté}, vérifié contre ses propres déclarations annuelles.}
\pts{
\item Trois compteurs, trois regards : \textbf{408} membres tradent des actions ; \textbf{243} partent
d'une photo de départ (les autres, de zéro) ; \textbf{376} ont un portefeuille copiable au bout du compte.
\item \textbf{129\,538} transactions les font vivre jour après jour.
\item L'Acte III pose une seule question : \textbf{les copier rapporte-t-il ?} On garde exactement la même
donnée, on change seulement ce qu'on en fait.
}
\end{tcolorbox}}
\vspace{2.6mm}

% ═══════════════ ACTE III — LA STRATÉGIE ═══════════════
\acte{ACTE III --- la stratégie : on empile des tests, une mesure à chaque fois}{orange}
\rangee{\begin{etapebox}{8}{orange}{8 $\cdot$ LA RÈGLE --- copier les mieux classés}
\reponse{Le portefeuille reconstruit d'un membre est un \textbf{fonds} : on en achète des parts. La règle
testée est la plus simple du projet.}
\phare{376}{membres dont le portefeuille peut être copié}{orange}
\vis{fig_10_valeur_part.png}
\pts{
\item \textbf{Éligible} à chaque 31 décembre si $\ge$ 126 jours d'historique et Sharpe $> 0{,}4$ --- sans
jamais regarder l'avenir.
\item Chaque année, trois gestes : \textbf{compter} le capital, le \textbf{redéployer} à parts égales sur
les mieux classés, \textbf{tenir} un an. Rien n'entre, rien ne sort.
}
\src{source : cellules \og un membre = un fonds \fg{} et \og le moteur en trois temps \fg{}.}
\end{etapebox}}{\begin{etapebox}{9}{orange}{9 $\cdot$ LE RÉGLAGE --- face au SPY}
\reponse{Le nombre de membres copiés est choisi chaque année sur le passé seulement. Voici ce que donne ce
réglage de base contre le \textbf{SPY}, l'indice des grandes actions américaines (dividendes réinvestis).}
\phare{$+0{,}2$ \%}{d'écart au SPY par an, en moyenne sur 11 années}{orange}
\vis{fig_10_nav.png}
\pts{
\item 100~\$ placés au départ donnent 451~\$ au bout de 10,5 ans, contre 433~\$ pour le SPY.
\item Quatre années au-dessus du marché, sept en dessous : la moyenne ne se lit pas sans son incertitude
(étape 11).
}
\src{source : cellule \og le réglage de base \fg{}.}
\end{etapebox}}
\vspace{2mm}
\rangee{\begin{etapebox}{10}{orange}{10 $\cdot$ COMBIEN COPIER --- le nombre K}
\reponse{On change un seul réglage : le nombre de membres copiés, fixé de 4 à 20, et on regarde ce que
devient l'écart au marché.}
\phare{K = 5}{le nombre qui produit le plus grand écart mesuré}{orange}
\vis{fig_10_balayage_k.png}
\pts{
\item L'écart est maximal au milieu de la grille : $K = 5$ donne $+1{,}3$ \%/an. Copier très peu ne fait
pas grandir l'écart, cela le rend plus instable ; copier beaucoup le dilue puis l'inverse.
\item Aucune valeur de $K$ n'approche le seuil de 1,81 : le meilleur cas atteint $t = 0{,}63$.
}
\src{source : cellule \og le balayage de K \fg{}.}
\end{etapebox}}{\begin{etapebox}{11}{orange}{11 $\cdot$ LA SOURCE, ET LA PORTÉE DU TEST}
\reponse{Même règle, quatre jeux de données (courbes ci-dessous). Puis la question qu'aucune mesure ne
règle : avec onze années, qu'aurait-on pu détecter ?}
\phare{4,6 \%}{l'écart annuel minimal détectable huit fois sur dix}{orange}
\vis{fig_10_sources.png}
{\footnotesize écart/an $\cdot$ $t$ : rapides $+0{,}9$/$0{,}37$ $\cdot$ $+$photo $-0{,}7$/$-0{,}23$ $\cdot$
$+$annuels $+0{,}1$/$0{,}07$ $\cdot$ complet $+0{,}2$/$0{,}13$ ($\beta \approx 0{,}9$ partout).}
\vspace{0.5mm}
\pts{
\item Les quatre sources restent \textbf{toutes proches de zéro} ($-0{,}7$ à $+0{,}9$ \%/an) : enrichir la
donnée change ce qu'on copie, pas l'écart au marché.
\item Le plus haut $t$ est $0{,}37$, loin de 1,81. Mais sous 4,6 \%/an, un écart réel serait resté
invisible : le $t$ faible \textbf{borne} ce qu'on voit, il ne prouve pas qu'il n'y a rien.
}
\src{source : cellules \og le réglage de base \fg{} (4 sources) et \og la puissance du test \fg{} ($\sigma = 5{,}7$ \%/an).}
\end{etapebox}}
\vspace{2mm}

% ═══════════════ CLÔTURE ═══════════════
\begin{tcolorbox}[enhanced, colback=rouge!3, colframe=rouge!45, boxrule=0.5pt, arc=2pt,
                  left=3mm, right=3mm, top=1.6mm, bottom=1.6mm,
                  colbacktitle=rouge, coltitle=white, fonttitle=\footnotesize\bfseries,
                  title={CE QUE CETTE PREMIÈRE VERSION NE DIT PAS}]
\begin{multicols}{2}
\pts{
\item \textbf{Datation à l'opération}, pas à la publication : les chiffres sont une \textbf{borne haute}.
\item \textbf{Aucun coût de transaction} n'est appliqué.
\item \textbf{Aucun placebo}, une seule pondération (parts égales), une seule règle (le Sharpe).
\item \textbf{Le juge reste la vérification} (étape 7) : elle dit si la reconstruction colle aux
déclarations, que la stratégie gagne ou non.
}
\end{multicols}
\vspace{-1mm}
{\scriptsize\color{grisf}\textbf{Pourquoi croire ces chiffres} : empreinte sha256 des tables conforme au
manifeste $\cdot$ cascades et tableaux vérifiés par assert $\cdot$ classe d'actif décidée avant tout prix
$\cdot$ moteur testé de trois façons ($r=0$ à l'injection, recalcul naïf, photo recomptée) $\cdot$ aucun
chiffre écrit en dur : tout se régénère en ré-exécutant le notebook.}
\end{tcolorbox}

\end{document}
